\documentclass[a4paper,10.5pt]{article}
\usepackage[utf8]{inputenc}
\usepackage[T1]{fontenc}
\usepackage[french]{babel}
\usepackage[left=1cm,right=1cm,top=.5cm,bottom=0cm]{geometry}
\usepackage{enumitem}
\usepackage{hyperref}
\usepackage{xcolor}
\usepackage{titlesec}
\usepackage{fontawesome5}
\usepackage{lmodern}
\usepackage[normalem]{ulem}

% --- Couleurs Siemens (teal corporate) ---
\definecolor{primary}{RGB}{0, 153, 153}
\definecolor{secondary}{RGB}{80, 80, 80}

% --- Config Liens ---
\hypersetup{colorlinks=true, linkcolor=primary, filecolor=primary, urlcolor=primary}

% --- Titres ---
\titleformat{\section}{\large\bfseries\color{primary}\uppercase}{}{0em}{}[\titlerule]
\titlespacing{\section}{0pt}{8pt}{5pt}

% --- Commandes ---
\newcommand{\resumeItem}[1]{\item\small{#1}}

\begin{document}

% --- EN-TÊTE ---
\begin{center}
    {\Large \textbf{Loïc Aron MBASSI EWOLO}} \\ \vspace{4pt}
    \textbf{\large Stage Systèmes Embarqués -- Capteurs \& Protocoles SPI} \\
    \textit{\small Stage disponible dès avril 2026 pour 4 à 6 mois} \\ \vspace{4pt}
    \small
    \faMapMarker\ Cergy, France (Mobile IDF) \quad
    \faPhone\ (+33) 7 59 52 33 98 \quad
    \faEnvelope\ \href{mailto:loicaron.mbassiewolo@ensea.fr}{loicaron.mbassiewolo@ensea.fr} \\ \vspace{2pt}
    \faLinkedin\ \href{https://www.linkedin.com/in/loic-mbassi/}{linkedin.com/in/loic-mbassi} \quad
    \faGithub\ \href{https://github.com/Nameless0l}{github.com/Nameless0l} \quad
    \faGlobe\ \href{https://mbassiloic.tech}{mbassiloic.tech}
\end{center}

\vspace{-6pt}

% --- PROFIL ---
\noindent\small{Drivers SPI en C sur STM32, acquisition de capteurs et validation sur prototype -- compétences acquises en laboratoire embarqué et en compétition robotique. Double diplôme ENSEA / Polytechnique Yaoundé, spécialité Systèmes Embarqués. Je suis disponible en stage dès avril 2026 pour rejoindre le département R\&D de \textbf{Siemens} à Vélizy.}

% --- FORMATION ---
\section{Formation}
\begin{itemize}[leftmargin=*, label=\tiny$\bullet$, nosep]
    \item \textbf{ENSEA -- École Nationale Supérieure de l'Électronique et de ses Applications} \hfill \textit{2025 -- 2027} \\
    \small{Double diplôme d'Ingénieur -- Systèmes Embarqués, Signal, IA. \textit{Cursus :} Architecture microcontrôleurs ARM Cortex-M, Protocoles (SPI, I2C, UART), FPGA/VHDL, Traitement du signal, Électronique analogique/numérique.}
    \item \textbf{École Polytechnique de Yaoundé (1\textsuperscript{ère} école d'ingénieur CEMAC)} \hfill \textit{2021 -- 2025} \\
    \small{Diplôme d'Ingénieur de Conception (Master 2) : Mathématiques Appliquées, Génie Logiciel, Algorithmique avancée.}
\end{itemize}

% --- COMPÉTENCES ---
\section{Compétences Techniques}
\begin{itemize}[leftmargin=*, nosep]
    \item \small{\textbf{Langages :} C (firmware embarqué), C++, Python}
    \item \small{\textbf{Microcontrôleurs :} STM32, Arduino, Raspberry Pi, Nvidia Jetson -- Renesas RX231 (à prendre en main)}
    \item \small{\textbf{Protocoles \& Interfaces :} SPI, I2C, UART, GPIO, PWM, interruptions, capteurs IMU, flex sensors}
    \item \small{\textbf{Outils :} GCC, Makefile, Git, Linux, GDB, oscilloscope (notions mesure en labo)}
    \item \small{\textbf{Visualisation \& Tests :} OpenGL (visualisation temps réel), tests unitaires, validation sur cible, comparaison cahier des charges}
    \item \small{\textbf{Langues :} Français (natif), Anglais (professionnel/technique)}
\end{itemize}

% --- PROJETS EMBARQUÉS ---
\section{Projets Embarqués}
\begin{itemize}[leftmargin=*, label={-}]

\item \textbf{Harmony Gloves -- Firmware \& Protocoles SPI/I2C (Gants Instrumentés)} \hfill \textit{2024} \\
\textit{Laboratoire IoT -- ENSPY} \hfill \textit{STM32/Arduino, C, SPI/I2C, Capteurs IMU/Flex, Soudure}
\vspace{-2pt}
    \begin{itemize}[leftmargin=0.4cm, nosep, label=\tiny$\bullet$]
        \item \small{Participation à l'écriture des drivers \textbf{SPI/I2C} en C sur STM32 pour l'acquisition de capteurs IMU et flex sensors.}
        \item \small{Gestion des interruptions, timing et synchronisation entre microcontrôleur et périphériques.}
        \item \small{Cycle complet : soudure composants, câblage, débogage bas niveau à l'oscilloscope, validation fonctionnelle.}
    \end{itemize}
\vspace{3pt}

\item \textbf{Stage Recherche IoT Lab -- Microcontrôleurs \& Contraintes Temps Réel} \hfill \textit{Oct. 2023 -- Jan. 2024} \\
\textit{ENSPY -- Laboratoire IoT} \hfill \textit{C, Arduino, Raspberry Pi, TensorFlow Lite, Python}
\vspace{-2pt}
    \begin{itemize}[leftmargin=0.4cm, nosep, label=\tiny$\bullet$]
        \item \small{Portage et optimisation de code C pour microcontrôleurs à ressources limitées (RAM, flash, fréquence CPU).}
        \item \small{Vérification des contraintes temps réel sur cible -- comparaison performances mesurées vs cahier des charges.}
    \end{itemize}
\vspace{3pt}

\item \textbf{Upgrade Jetson -- Challenge CoHoMa (Robotique Autonome Militaire)} \hfill \textit{2025 -- En cours} \\
\textit{ENSEA / Armée Française} \hfill \textit{C/C++, Nvidia Jetson, Lidar 3D, Caméra profondeur, Docker}
\vspace{-2pt}
    \begin{itemize}[leftmargin=0.4cm, nosep, label=\tiny$\bullet$]
        \item \small{Développement firmware C/C++ pour intégration de capteurs Lidar 3D (VLP16) et caméra de profondeur.}
        \item \small{Tests d'intégration hardware/software en conditions opérationnelles -- validation comportement sur prototype réel.}
    \end{itemize}
\vspace{3pt}

\item \textbf{PROJET-TAXI -- Simulation Multiprocessus C avec Visualisation} \hfill \textit{2024} \\
\textit{Projet Académique -- Polytechnique Yaoundé} \hfill \textit{C, Linux, IPC (mémoire partagée, signaux), OpenGL}
\vspace{-2pt}
    \begin{itemize}[leftmargin=0.4cm, nosep, label=\tiny$\bullet$]
        \item \small{Développement d'une interface de visualisation temps réel en \textbf{OpenGL} couplée à un backend C pur multiprocessus.}
        \item \small{Gestion mémoire partagée, signaux POSIX et tests de robustesse sous charge.}
    \end{itemize}

\end{itemize}

% --- PROJET COMPLÉMENTAIRE ---
\section{Projet Académique Complémentaire}
\begin{itemize}[leftmargin=*, label={-}]

\item \textbf{Fault Detection \& Fault-Tolerant Control -- Drone} \hfill \textit{2025} \\
\textit{Projet Option -- ENSEA} \hfill \textit{MATLAB/Simulink, Contrôle robuste, FDI/FTC}
\vspace{-2pt}
    \begin{itemize}[leftmargin=0.4cm, nosep, label=\tiny$\bullet$]
        \item \small{Modélisation dynamique, détection de défauts capteurs/actionneurs (FDI), contrôle tolérant aux pannes.}
        \item \small{Simulation de scénarios de défaillance -- évaluation comparative avec les spécifications initiales.}
    \end{itemize}

\end{itemize}

% --- DISTINCTIONS ---
\section{Distinctions \& Leadership}
\begin{itemize}[leftmargin=*, label=\tiny$\bullet$, nosep]
    \item \textbf{1\textsuperscript{ère} Place} -- IndabaX Africa 2025 (Hackathon IA Santé) \hfill \textit{2025}
    \item \textbf{1\textsuperscript{ère} Place} -- CITS National Hackathon (Smart City, contre 75 équipes) \hfill \textit{2024}
    \item \textbf{Lead AI Cell} -- Polytechnique Yaoundé : animation workshops C/C++, Git, Architecture logicielle \hfill \textit{2023-2024}
\end{itemize}

\end{document}
